\section{What It Means, and What Is Still Open}\label{sec:discussion}

\paragraph{What the offload makes possible.}
On two 96\,GB GPUs neither the original checkpoint nor the \dv{cfg.exl3_bits}-bit pack can keep all its experts in GPU
memory, and moving about a third of the experts to the host closes the last \dv{expl.noplan_over_per_gpu_gib}\,GiB per
card. The result serves a model whose unquantized version ranks first on agentic coding on LiveBench, on two GPUs that
can be rented by the hour. On the subsets we tested it matches the vendor's service within the uncertainty, and a long
context stays interactive once it is cached.

\paragraph{Practical guidance.}
For someone deploying this stack, four rules follow from the measurements. Each comes from a single boot on a specific
pod, so it is a starting point rather than a guarantee, and each names the check or the table to look at.
\begin{enumerate}
\item \emph{Check where the two GPUs sit before measuring or serving.} The same configuration had a median time per
token \dv{ctl.pod.cross.A.N1.tpot_ratio} to \dv{ctl.pod.cross.A.N32.tpot_ratio} times higher on a host where
\texttt{nvidia-smi topo -m} reports \texttt{\dv{ctl.pod.s2.gpu_link}} between the GPUs instead of
\texttt{\dv{ctl.pod.s1.gpu_link}} (and whose CPU also differed). Run \texttt{nvidia-smi topo -m} on the pod and
read the cell between GPU0 and GPU1. Our runners record it with the rest of the host (\datafile{experiments/lib.sh},
Tables~\ref{tab:pod-hardware} and~\ref{tab:cross-pod}).
\item \emph{Keep long contexts cached.} Once a 128K or 256K-token context is in the prefix cache, four streams keep a
median of \dv{agent.E13.own128k.prov.rate_p50} and \dv{agent.E13.cached256k.prov.rate_p50} tokens/s with TTFT p95
under two seconds, while a fresh 128K-token prompt costs \dv{ctl.E12.long.prefill-long-131072.closed:1.ttft_p50}\,s
of reading before its first token. Route the requests of one repository to the same server so that its prefix stays
in the cache (Figure~\ref{fig:read-write}, Tables~\ref{tab:ctl-E12} and~\ref{tab:agent-levels}).
\item \emph{Size interactive capacity in simultaneous streams, not tokens per month.} About
\dv{agent.E13.agent.Cprov} streams at 40 tokens/s per pod for agent turns, which under the observed usage profile
serve about \dv{agent.E13.agent.seats_conc_e1} intermittently active developers. Measure the team's active hours
and read the matching row of Table~\ref{tab:agent-duty}, and size the pods with Table~\ref{tab:agent-cost}.
\item \emph{Do not expect upstream vLLM's zero-copy offload to replace the recipe's hooks} on these pins. It does not
serve the model at all (Table~\ref{tab:ctl-E3}, with both KV block sizes).
\end{enumerate}

\paragraph{What we cannot yet explain.}
Each stream slows down as users are added, and we do not know which resource is responsible. The candidates are the
reads of cold experts over PCIe, the growing batch that speculative decoding must verify, interference from reading
new prompts, and the eager execution of decode steps larger than the biggest CUDA graph. Memory capacity is not among
them, since the pools are preallocated and did not move with load, and the KV cache was never the limit. The
experiments designed to separate the candidates (DSpark off, other placement policies, another SM budget for the
host-expert stream) did not run (Section~\ref{sec:not-run}). The same is true of placement. The plans are static and
derived from usage profiles, the recipe's own notes report that a plan tuned on one kind of traffic does worse on
another~\citep{diffbot2026dsv41exl3}, and recent systems argue that static placement is suboptimal in
general~\citep{he2026expertflow,zhong2025hybrimoe,zhu2026dali}. Our two nested plans show that a larger share on the
host costs throughput, and they say nothing about how much a better choice of experts would recover.
Section~\ref{sec:limits} lists every question the campaign leaves open, and Table~\ref{tab:plan} the experiments
written to answer them.

\section{Conclusion}\label{sec:conclusion}

A \dv{hub.orig.params_b}B-parameter model does not fit on two 96\,GB GPUs, even with its experts at
\dv{cfg.exl3_bits} bits, and a stack built by others makes it fit by moving its least-used experts to a back room in
host memory. The configuration boots and serves, matches the vendor's service within the uncertainty on the subsets we
tested, and for a single request keeps writing at nearly the same pace after reading a 256K-token prompt. It pays in
reading long prompts, in the speed of each stream as users are added, and in a sensitivity to the host as large as
most settings we varied. For coding agents, one server pours enough for a large team but has four stools, which
supports about \dv{agent.E13.agent.seats_conc_e1} developers who use the model now and then, close to the price of
their subscriptions, with larger pooled deployments below it. All of this rests on one boot per experiment.

The question in the introduction asked at what cost, and one part of the answer is still missing: how much speed the
offload costs against the same model held entirely in GPU memory. Measuring it needs larger GPUs and the original
checkpoint. The protocol is written, the variant is scripted to build from pinned sources, and the workload and the
analysis are the ones used throughout this paper. It has not yet been run, and on a prepared pod with two B300 GPUs
it is one command, \texttt{make bench-E9}.
